\section*{编译原理笔记}

\subsection*{语言的文法}

\subsection*{短语到句柄}
\textbf{语法树 最左推导 最右推导 短语 直接短语 素语句柄} \par
最左推导: 见左拆左。 \par
最右推导: 见右拆右。 \par
句型: 一棵语法树对应的整串符号。 \par
短语: 任意一棵子树的全部叶子节点串。 \par
直接短语: 子树只有父子两层(层深度只有 2), 并且子层为叶子节点。 \par
素短语: 子树至少包含一个终结符, 并且只有一个短语, 其内部不再包含更小的素短语。 \par
句柄: 句型最左边那个直接短语。 \par
\par
例：\par
文法 \( G[E]: E \to E+E \mid E*E \mid (E) \mid i \), 字符串: i+i*i \par
最左推导: \( E \to E+E \to i+E \to i+E*E \to i+i*E \to i+i*i \) \par
最右推导: \( E \to E+E \to E+E*E \to E+E*i \to E+i*i \to i+i*i \) \par

\begin{tikzpicture}[level distance=1.2cm,sibling distance=1cm]
% 左语法树：E→E+E, 右E→E*E
\node{E}
child{node{E} child{node{i}}}
child{node{$+$}}
child{node{E}
    child{node{E} child{node{i}}}
    child{node{$*$}}
    child{node{E} child{node{i}}}
};
\end{tikzpicture}
\hspace{3cm}
\begin{tikzpicture}[level distance=1.2cm,sibling distance=1cm]
% 右语法树：E→E+E, 右E→E*E
\node{E}
child{node{E} child{node{i}}}
child{node{$+$}}
child{node{E}
    child{node{E} child{node{i}}}
    child{node{$*$}}
    child{node{E} child{node{i}}}
};
\end{tikzpicture}

\todo[inline]{注, 如果最左推导和最右推导画出来两个语法树不同, 则有二义性, 相同则无二义性, 如上图所示。}

短语: i+i*i, i*i, i \par
直接短语：i \par
素短语: i \par
句柄: i \par

\subsubsection*{消除左递归}

\( A \to A\alpha \mid \beta \quad \text{化为}
\begin{cases}
A \to \beta A' \\
A' \to \alpha A' \mid \varepsilon
\end{cases} \)
\hspace{3cm}
\( (\alpha,\beta \text{为任意字符串}) \)
\todo[inline]{产生式有两条选择: 一个生成 \( \beta \), 一个生成 \( \alpha \), \( \alpha \) 不断进行循环, 生成无数个 \( \alpha \), 因此最终生成 \( \beta \dot{\alpha} \)} \par

\[
\begin{aligned}
\text{例1：}G(E)\quad
\begin{cases}
E\to E+T\mid T\\
T\to T*F\mid F\\
F\to (E)\mid i
\end{cases}
\Rightarrow
\begin{cases}
E\to TE' \\
E'\to +TE'\mid \varepsilon \\
T\to FT' \\
T'\to *FT'\mid \varepsilon \\
F\to (E)\mid i
\end{cases}
\end{aligned}
\]

\subsubsection*{提取公因式}

\[
A \to \delta\beta_1 \mid \delta\beta_2 \mid \dots \mid \delta\beta_i \mid \gamma_1 \mid \gamma_2 \mid \dots \mid \gamma_j
\quad \text{化为}
\begin{cases}
A \to \delta A' \mid \gamma_1 \mid \gamma_2 \mid \dots \mid \gamma_j \\
A' \to \beta_1 \mid \beta_2 \mid \dots \mid \beta_i
\end{cases}
\]


\[
\begin{aligned}
&\text{例1：已知文法} \\
&G(A):
\begin{cases}
A \to aAB \mid a \\
B \to Bb \mid d
\end{cases}
\quad\Downarrow\\[6pt]
&\begin{aligned}
A &\to aA' \quad &A' \to AB \mid \varepsilon \\
B &\to dB' \quad &B' \to bB' \mid \varepsilon
\end{aligned}
\end{aligned}
\]

\[
\begin{aligned}
&\text{例2：}\\
&G[A]:\ A \to aAB \mid a \mid b
\quad\Downarrow\\[4pt]
&\begin{aligned}
A &\to aA' \mid b \\
A' &\to AB \mid \varepsilon
\end{aligned}
\end{aligned}
\]

\subsection*{FIRST FOLLOW}
First集: 看产生式左边。 \par
Follow集: 看产生式右边。 \par


\subsection*{DFA 转换}

\subsection*{LL(1)}

\subsection*{LR(0)}

\subsection*{SLR(1)}

\subsection*{LR(1)}

\subsection*{语法制导的翻译与属性文法}

\subsection*{S/L 属性文法与翻译模式}

\subsection*{源程序的中间表示}

\subsection*{四元式}

\subsection*{算符优先}

\subsection*{波兰表达式}
